% Imporditud: tools/import_eksperimendid.py
% Allikas: Eesti füüsikaolümpiaadi eksperimendi ülesannete
%   kogu (Rael Kalda, 2023), ülesanne Ü32, lahendus L32.
\setAuthor{Erkki Tempel}
\setRound{lõppvoor}
\setYear{2022}
\setNumber{P E2}
\setDifficulty{1}
\setTopic{vedelike mehaanika}
\setVahendid{Tundmatu keha, niit, kaal (täpsus $\pm_0$,1 g), plastikust tops veega (200 ml)}
\setPunktid{12}
\setAllikas{Eksperimendi kogu Ü32, lahendus lk 46}
\setLahendusseis{toores}

\prob{Keha tihedus}
Määrake tundmatu keha tihedus.
\textit{NB:} Joonlaua kasutamine ei ole lubatud.

\hint

\solu
Tiheduse saame leida valemiga $\rho$ = m , kus m on keha mass ning V keha ruumala. V Mõõdame kaaluga tundmatu keha massi m.

Seejärel mõõdame kaaluga vee ja topsi kogumassi M . Kinnitame tundmatu keha

niidi külge ning riputame keha vette hõljuma (keha ripub niidi küljes ning anuma

põhja ei puuduta). Mõõdame nüüd näiva vee (veetase topsis tõusis) ja topsi kogu-

massi M $'$.

Kui keha on vette rippuma pandud, siis mõjub kaalule jõud F $'$ = M g + $\rho$vV g, kus V on keha ruumala, seega M $'$ = M + $\rho$vV.

Avaldame viimasest seosest keha ruumala

M$' -$ M V=

$\rho$v

Teades nüüd keha massi m ja ruumala V , saame leida keha tiheduse

$\rho$ = m = m$\rho$v V M$' -$ M
\probend
